\section{Analysis}

\begin{frame}{Planned Bayesian Hierarchical Model}
	\small
	For respondent $i$ in group $g(i)$, and outcome $k \in \{\text{H1N1},\text{Seasonal}\}$:
	$$
		y_{ik} \sim \mathrm{Bernoulli}(p_{ik}), \qquad
		\mathrm{logit}(p_{ik}) = \alpha_k + \mathbf{x}_i^\top \boldsymbol{\beta}_k + u_{g(i),k}
	$$
	$$
		u_{g,k} \sim \mathcal{N}(0,\sigma_{u,k}^2), \qquad
		\beta_{j,k} \sim \mathcal{N}(0,\tau_k^2)
	$$
	\begin{itemize}
		\item Partial pooling allows stable estimates for smaller subgroups.
		\item Candidate group structure: HHS region and selected demographic strata.
		\item Priors will be weakly informative to regularize high-dimensional predictors.
	\end{itemize}
\end{frame}

\begin{frame}{Inference and Model-Checking Plan}
	\begin{itemize}
		\item Inference approach: MCMC posterior sampling (NUTS) for both outcomes.
		\item Planned diagnostics:
		\begin{itemize}
			\item Convergence checks with $\hat{R}$ near 1.00 and sufficient effective sample size.
			\item Posterior predictive checks for calibration and distributional fit.
			\item Sensitivity analysis under alternative prior scales.
		\end{itemize}
		\item Outputs to interpret:
		\begin{itemize}
			\item Posterior means/intervals for key predictors.
			\item Group-level random effect uncertainty.
			\item Predicted vaccination probabilities with credible intervals.
		\end{itemize}
	\end{itemize}
\end{frame}

\begin{frame}{Planned Baselines and Comparison Criteria}
	\begin{itemize}
		\item Baseline 1: regularized logistic regression.
		\item Baseline 2: random forest classifier.
		\item Proposed comparison dimensions:
		\begin{itemize}
			\item Predictive performance: macro ROC AUC across both vaccine targets.
			\item Calibration quality of predicted probabilities.
			\item Interpretability and uncertainty quantification for public health use.
		\end{itemize}
		\item Hypothesis: hierarchical Bayesian model will offer stronger uncertainty-aware interpretation with competitive predictive performance.
	\end{itemize}
\end{frame}